\chapter{Static assets and frontend integration}

\noindent\textbf{Learning path: Fast path: browser branch.} Read now when adding static assets or progressive browser enhancement.\par\medskip

A web application produces more than dynamic HTML and JSON responses. CSS, icons, images, fonts, source maps, or artifacts from a separate frontend build have a different lifecycle: they arrive as part of a release, ideally remain immutable, and can be cached aggressively by the browser.

Velran therefore serves static assets from a separate, \textbf{read-only static root}. This is not the same area as the writable AppFs data root, and it is not the same concept as user-uploaded media.

\noindent\textbf{First read:} learn the three file lifecycles, fingerprinted caching, and the V1 JavaScript/CSP boundary. Server-vs-proxy serving examples are reference material for deployment.

\section{Three separate file lifecycles}

In production, keep at least three file categories separate:

\begin{center}
\begin{tabularx}{\textwidth}{@{}lYYY@{}}
\toprule
 & Example & Writable? & Lifecycle \\
\midrule
release & CSS, font, logo & read-only & replaced on deploy \\
runtime data & export, generated file & according to AppFs policy & changes while the app runs \\
user media & uploaded image & controlled writable & lives with business data \\
\bottomrule
\end{tabularx}
\end{center}

A typical directory structure is:

\begin{lstlisting}
/usr/local/bin/velran-server
/usr/local/bin/velran-cli
/srv/myapp/releases/2026-09-04/
  main.vrn
  migrations/
  public/
    css/
    img/
    fonts/
/srv/myapp/data/
/run/secrets/velran/
\end{lstlisting}

Treat \code{public/} as a deployment artifact. \code{data/}, by contrast, is runtime state. If they are mixed, an upload may overwrite a release file, or a rollback may accidentally roll user data backward.

\section{Configuring the static root}

Static namespace configuration lives in a separate server-config section:

\begin{lstlisting}
[static_assets]
root = "/srv/myapp/current/public"
url_prefix = "/assets/"
max_asset_bytes = 8388608
max_age_secs = 300
immutable_max_age_secs = 31536000
precompressed = true
\end{lstlisting}

The same options are available through the CLI. Important flags include:

\begin{lstlisting}
--static-root
--static-url-prefix
--max-static-asset-bytes
--static-max-age-secs
--static-immutable-max-age-secs
\end{lstlisting}

The static prefix is a reserved route namespace. If the prefix is \code{/assets/}, the application must not declare a conflicting dynamic or static route. The server should fail closed on collision: startup failure is better than ambiguous route selection per request.

\section{Fingerprints: the key to long caching}

Tie static content and URL together. For example:

\begin{lstlisting}
public/
  css/app.01234567.css
  img/logo.a91f028c.svg
  fonts/ui-f1e2d3c4.woff2
\end{lstlisting}

Velran treats a filename as fingerprinted when it contains a token of at least eight hexadecimal characters delimited by a dot, hyphen, or underscore. Such a file may receive a long immutable cache policy:

\begin{lstlisting}
Cache-Control: public, max-age=31536000, immutable
\end{lstlisting}

For a non-fingerprinted file the safer default is short and revalidatable:

\begin{lstlisting}
Cache-Control: public, max-age=300, must-revalidate
\end{lstlisting}

The key rule is: \textbf{do not replace content behind the same URL while using a one-year cache}. A new CSS build should get a new fingerprinted filename. Old files may remain in the release long enough for old HTML representations to age out.

\section{A minimal asset build pipeline}

Velran does not prescribe a frontend build system. The pipeline only needs to produce deterministic output into \code{public/}.

For a small application, source might look like:

\begin{lstlisting}
assets-src/
  app.css
  logo.svg
public/
\end{lstlisting}

A release build may:

\begin{enumerate}
  \item create a clean \code{public/} directory;
  \item build or copy CSS/fonts/images;
  \item compute a content hash;
  \item create fingerprinted filenames;
  \item optionally produce Brotli/gzip siblings;
  \item use the fingerprinted URLs from the application.
\end{enumerate}

The hash is not a security token. Its purpose is to identify the content version in the filename.

\section{ETag and conditional GET}

Every static representation receives its own \code{ETag}. The browser can later send:

\begin{lstlisting}
If-None-Match: "velran-..."
\end{lstlisting}

If the representation has not changed, the server may return \code{304 Not Modified} without a body. HEAD is supported alongside GET.

A static request does not create a session and does not send a session cookie. This matters for both performance and security: downloading CSS or a font should not become a session-lifecycle event.

\section{Precompressed Brotli and gzip}

The Velran server does not compress static files during the request. A release pipeline may prebuild:

\begin{lstlisting}
app.01234567.css
app.01234567.css.br
app.01234567.css.gz
\end{lstlisting}

Based on \code{Accept-Encoding}, the server chooses Brotli, then gzip, then the original representation. A compressed response includes headers such as:

\begin{lstlisting}
Content-Encoding: br
Vary: Accept-Encoding
\end{lstlisting}

This avoids per-request compression CPU cost and does not expose dynamic compression work as part of the web request resource surface.

\section{MIME, nosniff, and path confinement}

The server maps known web extensions to MIME types and uses \code{application/octet-stream} for unknown extensions. Responses always include:

\begin{lstlisting}
X-Content-Type-Options: nosniff
\end{lstlisting}

Static paths are intentionally strict. Valid:

\begin{lstlisting}
/assets/css/app.01234567.css
\end{lstlisting}

Rejected forms include \code{..}, encoded traversal, backslashes, empty path components, and control characters. V1 does not allow percent-encoding in static paths; the build pipeline should use URL-safe filenames.

On Linux, static-root access follows the same confinement philosophy as AppFs: serving assets must not become arbitrary host-filesystem reads through symlinks or traversal.

\section{CSS, images, and fonts in Velran HTML}

The V1 CSP allows same-origin style, image, font, and media assets. A normal release may therefore contain:

\begin{lstlisting}
/assets/css/app.01234567.css
/assets/img/logo.a91f028c.svg
/assets/fonts/ui-f1e2d3c4.woff2
\end{lstlisting}

HTML should always reference the fingerprinted URL that belongs to the release. Templates must not construct asset URLs from user input; the asset name is build data, not request data.

\section{Important V1 boundary: JavaScript may be served, but not executed in Velran HTML}

The static server knows the MIME types for \code{.js} and \code{.mjs}, so those build artifacts can physically be served. This does \textbf{not} mean that HTML rendered by Velran may execute scripts.

In the V1 HTML security model:

\begin{itemize}
  \item CSP does not grant a \code{script-src} permission;
  \item the compiler rejects \code{<script>} markup;
  \item there is no inline-script or ``unsafe'' escape hatch.
\end{itemize}

This is an intentional security boundary. JavaScript \code{fetch} examples in the previous chapter therefore have two valid interpretations:

\begin{enumerate}
  \item a separate frontend application runs on its own origin and calls a Velran JSON API;
  \item the code illustrates browser-API client behavior, not that a script tag can be embedded in V1 Velran HTML.
\end{enumerate}

For a classic Velran server-rendered application, build around HTML + CSS + typed form/action. If a full JavaScript frontend is required, deploy it separately and use Velran as a typed JSON backend.

\section{Should Velran serve assets, or Apache/Nginx?}

Two deployment models work.

\subsection*{1. Velran-owned static namespace}

This should be the default. The reverse proxy sends every request to Velran, which handles the \code{/assets/} namespace itself:

\begin{lstlisting}
Internet
   |
Apache / Nginx :443
   |
127.0.0.1:8080 Velran
   |-- application routes
   `-- /assets/*
\end{lstlisting}

The advantage is that fingerprint, ETag, precompressed, nosniff, and path policy stay in one place. This is usually more than sufficient for small and medium applications.

\subsection*{2. The reverse proxy serves assets directly}

At higher static traffic volumes, Apache or Nginx may serve the read-only release \code{public/} directory directly. If you do this, \textbf{you must reproduce} the important contract yourself:

\begin{itemize}
  \item separate read-only static root;
  \item traversal/symlink-safe path handling;
  \item correct MIME + \code{nosniff};
  \item long immutable cache for fingerprinted files;
  \item short/revalidatable cache for non-fingerprinted files;
  \item correct \code{Vary} for Brotli/gzip;
  \item no namespace collision with application routes.
\end{itemize}

This is an operator optimization, not a requirement for application correctness.

\section{Nginx example for direct static serving}

A simple starting point:

\Needspace{15\baselineskip}
\begin{lstlisting}
server {
    listen 443 ssl;
    server_name app.example;

    location /assets/ {
        alias /srv/myapp/current/public/;
        try_files $uri =404;
        add_header X-Content-Type-Options nosniff always;
    }

    location / {
        proxy_pass http://127.0.0.1:8080;
        proxy_set_header Host $host;
        proxy_set_header X-Forwarded-Proto https;
    }
}
\end{lstlisting}

This is only a topology baseline. Production configuration must align cache policy, precompressed files, and symlink/release rules with the same immutable deployment model.

\section{Apache example for direct static serving}

The idea is the same with Apache:

\Needspace{15\baselineskip}
\begin{lstlisting}
Alias /assets/ /srv/myapp/current/public/

<Directory /srv/myapp/current/public/>
    Require all granted
    Options -Indexes
    AllowOverride None
    Header always set X-Content-Type-Options "nosniff"
</Directory>

ProxyPass        / http://127.0.0.1:8080/
ProxyPassReverse / http://127.0.0.1:8080/
\end{lstlisting}

If Apache serves assets directly, the \code{Alias} rule must take effect before the proxy catch-all, and cache headers must be configured explicitly.

\section{Deploy application code and assets as one release}

A common asset failure is deploying HTML that references the new CSS filename while the static root still points to the old release -- or vice versa.

Use this order:

\begin{lstlisting}
build application
-> build/fingerprint assets
-> optional .br/.gz
-> immutable release directory
-> preflight
-> atomic current switch
-> controlled restart/reload
\end{lstlisting}

Application code and the \code{public/} directory should be part of the same release artifact. Rollback then restores both together.

\section{Checking with curl}

After deployment, check at least:

\begin{lstlisting}
curl -I https://app.example/assets/css/app.01234567.css
curl -I -H 'Accept-Encoding: br' \
  https://app.example/assets/css/app.01234567.css
curl -I -H 'If-None-Match: "velran-..."' \
  https://app.example/assets/css/app.01234567.css
\end{lstlisting}

Verify \code{Content-Type}, \code{Cache-Control}, \code{ETag}, \code{X-Content-Type-Options}, and, for compressed representations, \code{Content-Encoding} and \code{Vary}.

\section{Practical decision rule}

\begin{center}
\begin{tabularx}{\textwidth}{@{}lY@{}}
\toprule
Situation & Recommended approach \\
\midrule
server-rendered Velran page & Velran static root + CSS/image/font \\
small/medium production & proxy forwards everything to Velran \\
high static traffic & Apache/Nginx direct static with equivalent policy \\
full JavaScript frontend & separate frontend deployment + Velran JSON API \\
user upload & AppFs/media pipeline, not static root \\
\bottomrule
\end{tabularx}
\end{center}

The most important rule is: \textbf{a static asset is a deployment artifact, not application data}. Preserving that boundary simplifies caching, rollback, backup, and security at the same time.

\section{Verified companion example}
The current static-asset companion is \path{examples/static-assets/app.vrn}. Keep release assets read-only and separate from AppFs/runtime state exactly as shown by the deployment model in this chapter.

\section{Practice: classify every frontend file}
\paragraph{Exercise mode: supported.} Use the supported method defined in the preface.

Take the files used by a page and classify each as immutable deployment asset, generated HTML, JSON response, or mutable application data. Confirm that only immutable deployment assets live under the static root.

\section{Common mistake: letting frontend convenience erase trust boundaries}
A JavaScript frontend does not change server-side authorization requirements. Client validation and hidden controls are usability features; all authoritative validation, authorization, and mutation rules remain on the server boundary.

\section{Running project checkpoint: Secure Notes}
Add a small stylesheet and, if useful, progressive JavaScript enhancement to the notes UI. Keep the application functional without trusting client-side checks, and keep user-generated content outside the static asset tree.
